\section*{Problemas}

\subsection*{Enunciados de los problemas}

% Recordar
\begin{problema}
	\label{prob:3f1ec07}
	Enuncie, sin realizar ningún cálculo, la fórmula del estadístico de bondad de ajuste de Pearson $\chi^2$, y la regla general para determinar sus grados de libertad cuando el modelo teórico está completamente especificado de antemano (sin parámetros estimados de la muestra).
\end{problema}

% Comprender
\begin{problema}
	\label{prob:d4859df}
	Un científico de datos diseña una prueba de bondad de ajuste $\chi^2$ para un modelo de categorización discreta en cuatro clases ($k=4$, $\nu=3$) con $n=80$ observaciones, obteniendo frecuencias esperadas $E_1=45.2$, $E_2=30.5$, $E_3=3.2$, $E_4=1.1$. Explique por qué el cálculo directo del estadístico de Pearson viola las condiciones de regularidad asintótica (la Regla de Cochran), y proponga el procedimiento correctivo estándar, indicando cuántos grados de libertad tendría la prueba ajustada.
\end{problema}

% Aplicar
\begin{problema}
	\label{prob:da2eb15}
	Un genetista indaga si la distribución de grupos sanguíneos ABO en una cohorte sigue las proporciones de referencia nacional: 45\% A, 40\% B, 10\% AB, 5\% O. En una muestra de $n=200$ individuos se observan: A=95, B=75, AB=18, O=12. Realice la prueba de bondad de ajuste al nivel $\alpha=0.05$ (use $\chi^2_{0.05,3}=7.815$).
\end{problema}

% Analizar
\begin{problema}
	\label{prob:d6f0154}
	Un ingeniero registra el conteo diario de interrupciones de un servidor durante $n=100$ días: $0$ interrupciones (36 días), $1$ (40 días), $2$ (18 días), $\ge3$ (6 días, tómese como valor puntual $k=3$). Para verificar si la ocurrencia sigue una Poisson: (a) calcule $\hat\lambda$ por máxima verosimilitud; (b) calcule las frecuencias esperadas bajo Poisson($\hat\lambda$); (c) calcule $\chi^2$ y justifique por qué los grados de libertad correctos son $\nu=k-1-m$ con $m=1$ parámetro estimado. ¿Se acepta el modelo al $\alpha=0.05$?
\end{problema}

% Evaluar
\begin{problema}
	\label{prob:f622374}
	En tablas $2\times2$ con muestras pequeñas o medianas, el estadístico de Pearson $\chi^2$ (sin corrección) tiende a producir p-valores sistemáticamente más bajos que el valor exacto de Fisher, aumentando el riesgo de Error Tipo I. Evalúe por qué ocurre este sesgo, explicando la diferencia entre la naturaleza discreta y escalonada de la distribución hipergeométrica exacta y la naturaleza continua de la $\chi^2_1$ de referencia, y explique cómo la Corrección de Continuidad de Yates ($|ad-bc|-N/2$ en vez de $|ad-bc|$) corrige este sesgo.
\end{problema}

% Crear
\begin{problema}
	\label{prob:369a6a1}
	Diseñe un ejemplo original (contexto, categorías, frecuencias observadas hipotéticas) donde se aplique una prueba de bondad de ajuste $\chi^2$, distinto de los ejemplos de dados, tipos sanguíneos o servidores ya vistos en esta sección. Plantee $H_0$ y $H_a$, calcule las frecuencias esperadas y el estadístico $\chi^2$ para sus datos.
\end{problema}

\subsection*{Soluciones detalladas}

% Recordar
\begin{solproblema}[prob:3f1ec07]
	$\chi^2 = \sum_{i=1}^k \dfrac{(O_i-E_i)^2}{E_i}$. Cuando el modelo está completamente especificado (sin parámetros estimados), los grados de libertad son $\nu=k-1$, donde $k$ es el número de categorías.
\end{solproblema}

% Comprender
\begin{solproblema}[prob:d4859df]
	La aproximación continua de $\chi^2$ a la distribución discreta real de $Q=\sum(O_i-E_i)^2/E_i$ solo es confiable cuando ninguna $E_i$ es menor a $1.0$ y no más del $20\%$ de las celdas tienen $E_i<5$ (Regla de Cochran). Aquí, $E_3=3.2<5$ y $E_4=1.1<5$: el $50\%$ ($2/4$) de las celdas viola la cota, invalidando el uso directo de la tabla $\chi^2$. El procedimiento correctivo estándar es \textbf{agrupar categorías adyacentes de baja frecuencia}: fusionando las clases 3 y 4, $E_{3+4}=3.2+1.1=4.3$, resultando en $k'=3$ categorías efectivas y $\nu_{\text{corregido}}=k'-1=2$ grados de libertad.
\end{solproblema}

% Aplicar
\begin{solproblema}[prob:da2eb15]
	$E_A=90$, $E_B=80$, $E_{AB}=20$, $E_O=10$.
	\begin{align*}
		\chi^2 = \frac{(95-90)^2}{90}+\frac{(75-80)^2}{80}+\frac{(18-20)^2}{20}+\frac{(12-10)^2}{10} = 0.278+0.313+0.200+0.400 = 1.190.
	\end{align*}
	Como $1.190 < 7.815$, \textbf{no se rechaza} $H_0$: la muestra es consistente con las proporciones de referencia nacional.
\end{solproblema}

% Analizar
\begin{solproblema}[prob:d6f0154]
	(a) $\hat\lambda = \bar X = \dfrac{0(36)+1(40)+2(18)+3(6)}{100} = \dfrac{94}{100}=0.94$.
	(b) Con Poisson($0.94$): $E_0\approx39.06$, $E_1\approx36.72$, $E_2\approx17.26$, $E_{\ge3}\approx6.96$.
	(c) $\chi^2 = \dfrac{(36-39.06)^2}{39.06}+\dfrac{(40-36.72)^2}{36.72}+\dfrac{(18-17.26)^2}{17.26}+\dfrac{(6-6.96)^2}{6.96} \approx0.697$. Al estimar $m=1$ parámetro ($\hat\lambda$) de la propia muestra, cada parámetro estimado impone una ligadura adicional sobre las desviaciones $(O_k-E_k)$, reduciendo en una unidad extra la dimensión libre: $\nu=k-1-m=4-1-1=2$. Con $\chi^2_{0.05,2}=5.991$ y $0.697\ll5.991$, \textbf{se acepta} el modelo de Poisson.
\end{solproblema}

% Evaluar
\begin{solproblema}[prob:f622374]
	La distribución teórica exacta de una tabla $2\times2$ finita está dictada por la distribución Hipergeométrica de Fisher, que es \textbf{discreta y escalonada}. La $\chi^2_1$ de referencia, en cambio, es una curva \textbf{continua}. Al aproximar el salto escalonado discreto mediante esa curva continua, el área de la cola derecha (el p-valor) resulta sistemáticamente \textbf{más pequeña} que la probabilidad exacta real, produciendo p-valores artificialmente bajos y un exceso de falsos positivos. La Corrección de Yates ($|ad-bc|-N/2$ en vez de $|ad-bc|$) reduce la magnitud del numerador antes de elevarlo al cuadrado, disminuyendo el valor del estadístico y, por tanto, aumentando el p-valor calculado hacia el valor exacto de Fisher — corrigiendo el sesgo optimista de la aproximación continua sin corrección.
\end{solproblema}

% Crear
\begin{solproblema}[prob:369a6a1]
	\textbf{Ejemplo:} una tienda en línea quiere verificar si las devoluciones de productos se distribuyen uniformemente entre 4 categorías (Ropa, Electrónica, Hogar, Otros) según su política histórica de 40\%, 30\%, 20\%, 10\%. En $n=150$ devoluciones observadas: Ropa=68, Electrónica=40, Hogar=25, Otros=17. $H_0$: las proporciones son 40/30/20/10\%; $H_a$: al menos una difiere. Frecuencias esperadas: $E=[60,45,30,15]$.
	\begin{align*}
		\chi^2 = \frac{(68-60)^2}{60}+\frac{(40-45)^2}{45}+\frac{(25-30)^2}{30}+\frac{(17-15)^2}{15} \approx1.067+0.556+0.833+0.267 = 2.723.
	\end{align*}
	Con $\chi^2_{0.05,3}=7.815$, como $2.723<7.815$, no se rechaza $H_0$.
\end{solproblema}
